\thispagestyle{otFront}\gdef\otChapNum{0}\phantomsection\label{references}\pdfbookmark[0]{Références}{bm-references}
\otPageTitre{Références}{France · Seconde générale et technologique}
\otIntertitre{Le programme de référence}
Le livre s’appuie sur l’arrêté du 26 février 2026, publié au BO n° 14 du 2 avril 2026 (NOR MENE2602914A), applicable à la rentrée 2026-2027. Son annexe (pages 333 à 345 du BO) a été relue intégralement lors de l’audit du 23 octobre 2026. La matrice précédente repère les contenus, les capacités, les démonstrations, les exemples algorithmiques, les automatismes et les approfondissements dans le livre.\par\smallskip
\href{https://www.education.gouv.fr/bo/2026/Hebdo14/MENE2602914A}{Programme et annexe sur education.gouv.fr}\par\bigskip
\otIntertitre{Ressources pédagogiques consultées}
Les ressources Éduscol « Accompagner les élèves en mathématiques - 4e, 3e et 2de » et « Raisonner et démontrer en Seconde » ont servi à vérifier la variété des tâches et l’importance de l’argumentation. Les exercices de ce livre sont originaux ; ils ne reproduisent pas ces fiches.\par\smallskip
\href{https://eduscol.education.gouv.fr/6858/accompagner-les-eleves-en-mathematiques-4e-3e-et-2de}{Ressources Éduscol et fiches par domaines}\par\bigskip
\otIntertitre{Données réelles}
Deux exercices utilisent des données identifiées : Insee Focus n° 374, décembre 2025, pour les populations de Nantes et Bordeaux en 2023 ; dossier complet France, POP T4, édition du 27 août 2026, pour des croisements âge et sexe au RP2023. Le champ France hors Mayotte est précisé dans l’exercice concerné. Les autres enquêtes et modèles numériques sont construits à des fins pédagogiques.\par\smallskip
\href{https://www.insee.fr/fr/statistiques/8682718}{Insee Focus 374}\quad
\href{https://www.insee.fr/fr/statistiques/2011101?geo=FRANCE-1}{Dossier France, données 2023}\par\bigskip
\otIntertitre{Choix et vérifications}
Les calculs exacts, cas limites, énumérations et programmes ont été contrôlés avec les scripts fournis. Les preuves ont été relues ; les contrôles numériques finis ne remplacent pas une démonstration générale. Toutes les pages sont rendues et contrôlées avant livraison.\par\medskip
{\small\color{otMuted}Sources consultées le 21 octobre 2026 ; programme relu le 23 octobre 2026. Charte de la collection OptiTAB. Les deux bilans, leur barème et les séries de questions flash sont des choix pédagogiques.}\clearpage